\chapter{Manual de instalación y despliegue}
\label{app:manual_instalacion}

\ph{Este apéndice describe los pasos para instalar y desplegar el sistema en
un entorno limpio. Debe ser lo bastante detallado para que un tercero
(p.\,ej.\ otro estudiante o un evaluador) pueda reproducir el despliegue
sin contacto con el autor/a.}

\section{Requisitos previos}
\label{sec:requisitos_previos}

\ph{Lista del software, hardware y credenciales necesarios antes de empezar.
Indica versiones concretas siempre que sea posible.}

\ph{\begin{itemize}
  \item \textbf{Sistema operativo:} \ph{p.\,ej.\ Linux Ubuntu 24.04 LTS / macOS / Windows 11.}
  \item \textbf{Lenguaje y runtime:} \ph{p.\,ej.\ Python 3.x / Node 20 / Java 17.}
  \item \textbf{Gestor de paquetes:} \ph{pip / npm / maven\dots}
  \item \textbf{Servicios externos:} \ph{base de datos, broker de mensajería,
        almacenamiento\dots}
  \item \textbf{Otras herramientas:} \ph{Docker, Git, etc.}
\end{itemize}
}

\begin{itemize}
    \item \textbf{Sistema operativo:} el entorno utilizado durante el
    desarrollo y validación del sistema ha sido Ubuntu 24.04 LTS. También es
    posible utilizar otros sistemas compatibles con Docker.

    \item \textbf{Lenguaje y runtime:} no es necesario instalar manualmente
    Python ni sus dependencias en el sistema anfitrión, ya que el entorno de
    ejecución se encuentra definido dentro de los contenedores del proyecto.

    \item \textbf{Gestor de paquetes:} las dependencias necesarias para
    ejecutar la aplicación y el generador se instalan automáticamente durante
    la construcción de las imágenes Docker.

    \item \textbf{Servicios externos:} no es necesario instalar manualmente
    servicios adicionales. La configuración Docker del proyecto incluye los
    contenedores necesarios para los servicios utilizados por UVLHub, entre
    ellos la base de datos MariaDB, Redis, Elasticsearch, MailHog y Selenium.

    \item \textbf{Otras herramientas:} es necesario disponer de Git para
    obtener el código fuente y Docker junto con Docker Compose para construir y
    ejecutar el entorno de desarrollo.
\end{itemize}

\section{Obtención del código fuente}
\label{sec:obtencion_codigo}

\ph{Indica el repositorio donde está alojado el código y cómo clonarlo:}
En el siguiente enlace se puede encontra el repositorio del proyecto: https://github.com/antjimort/uvlhub

\begin{verbatim}
git clone git@github.com:antjimort/uvlhub.git
cd uvlhub
\end{verbatim}

\section{Instalación de dependencias y ejecución}
\label{sec:instalacion_dependencias}

\ph{Comandos para instalar las dependencias del proyecto en un entorno
aislado (virtualenv, contenedor, etc.). Si hay variables de entorno o
ficheros de configuración que el usuario debe rellenar, lístalos aquí.}

Las dependencias del proyecto se gestionan mediante Docker, por lo que no es
necesaria una instalación manual de paquetes. La construcción de los
contenedores instala automáticamente las dependencias definidas en la
configuración del proyecto.

Este comando permite construir las imágenes utilizando el archivo de configuración correspondiente que contiene todos los subcomandos internos que lo prepara todo. Posteriormente, levanta los contenedores.

\begin{verbatim}
docker compose -f docker/docker-compose.dev.yml up --build
\end{verbatim}

Para detener y eliminar los contenedores se usa el siguiente comando:

\begin{verbatim}
docker compose -f docker/docker-compose.dev.yml down
\end{verbatim}

\section{Despliegue en producción}
\label{sec:despliegue_produccion}

\ph{Si el sistema se despliega en un entorno productivo (servidor propio,
nube, PaaS\dots), describe el proceso paso a paso: construcción de la
imagen, publicación en el registro, despliegue en el host, configuración
del proxy/HTTPS, etc.}

La herramienta utilizada para el despliegue de la app principal ha sido Render, pero ha sido usada únicamente para las revisiones del tutor. Tambien se ha usado Redis, utilizado como almacenamiento externo para la gestión de las
sesiones de Flask para conservar el estado
del proceso de configuración del generador.

No obstante, la plataforma utilizada no constituye un requisito del sistema,
ya que la aplicación puede ser desplegada mediante otras alternativas
compatibles con contenedores.

El despliegue oficial de UVLHub es gestionado por DiversoLab.


\section{Verificación de la instalación}
\label{sec:verificacion}

\ph{Cómo comprobar que la instalación ha sido correcta: ejecución de la
suite de pruebas, acceso a una pantalla concreta, llamada a un endpoint
de \textit{health check}, etc.}

\section{Verificación de la instalación}
\label{sec:verificacion-instalacion}

Una vez levantados los contenedores mediante el comando indicado en la
Sección~\ref{sec:instalacion_dependencias}, se puede comprobar que la
instalación se ha realizado correctamente accediendo a la dirección local
proporcionada por la aplicación durante el arranque.

La consola muestra la dirección en la que se encuentra disponible el servicio
web, normalmente mediante la dirección
\texttt{http://172.18.0.10:5000} (la segunda que aparece). Se puede acceder a ella mediante cualquier navegador. Un ejemplo de la salida generada durante el
inicio de la aplicación se muestra en la Figura~\ref{fig:inicio-contenedor}.

\begin{figure}[H]
    \centering
    \includegraphics[width=0.9\textwidth]{salida_contenedor.png}
    \caption{Salida de consola tras el arranque correcto del contenedor web. Elaboración propia}
    \label{fig:inicio-contenedor}
\end{figure}